\section{Erd\H{o}s Problem \#884 --- continuation}

\section{ROUND-2 OBJECTIVE}

I continue on the proof side. Round-(round\_no-1) already established the notation and isolated the global obstruction: the general conjecture is still open unconditionally, while Tao's 2025 note gives only a conditional disproof under a qualitative Hardy--Littlewood prime tuples conjecture.

The most promising strict advance is therefore to enlarge the unconditional positive family. In this round I prove a new theorem that strictly extends the previous \(r^a s\) case.

\medskip

\noindent\textbf{Main new theorem.}
\emph{Let \(r\) be a prime, let \(a\geq 0\), and let \(m\geq 1\) be an integer with \((r,m)=1\). Write}
\[
M:=\tau(m).
\]
\emph{Let \(1=d_1<\cdots<d_t=n\) be the increasing list of divisors of}
\[
n=r^a m.
\]
\emph{Then}
\[
\sum_{1\leq i<j\leq t}\frac{1}{d_j-d_i}
\leq
H_{M-1}\sum_{1\leq i<t}\frac{1}{d_{i+1}-d_i}
+\frac{10}{3}M^2,
\]
\emph{where \(H_{M-1}:=\sum_{u=1}^{M-1}\frac1u\), with the convention \(H_0:=0\).}

Consequently, for every fixed \(M\), all integers of the form \(n=r^a m\) with \((r,m)=1\) and \(\tau(m)\leq M\) satisfy
\[
S_{\mathrm{all}}(n)\leq C_M\bigl(1+S_{\mathrm{gap}}(n)\bigr)
\]
with
\[
C_M:=H_{M-1}+\frac{10}{3}M^2.
\]

As a sharper special case, if \(n=p^a q^b\) with distinct primes \(p,q\) and
\[
m:=\min(a,b),
\]
then
\[
S_{\mathrm{all}}(n)\leq H_m\,S_{\mathrm{gap}}(n)+5(m+1),
\]
hence
\[
S_{\mathrm{all}}(n)\leq \bigl(H_m+5(m+1)\bigr)\bigl(1+S_{\mathrm{gap}}(n)\bigr).
\]

This strictly strengthens the previous round, which only handled the case \(m=1\).

\section{Round-(round\_no-1) FOUNDATION USED}

I use the following Round-(round\_no-1) results as fixed input.

\begin{enumerate}
\item The formal definitions
\[
S_{\mathrm{all}}(n):=\sum_{1\le i<j\le t}\frac{1}{d_j-d_i},
\qquad
S_{\mathrm{gap}}(n):=\sum_{1\le i<t}\frac{1}{d_{i+1}-d_i}.
\]
\item The previously isolated structural idea that the proof for \(n=r^a s\) works because the divisor set is the merge of two geometric progressions.
\item The previously established global status: the unconditional problem remains open, while Tao's 2025 note provides only a conditional disproof under a qualitative prime tuples conjecture.
\end{enumerate}

I do \emph{not} re-prove those facts here.

\section{NEW INSIGHT / TOOL (ROUND-2)}

The new ingredient is a general ``\(M\)-progressions'' lemma.

If the divisors of \(n\) can be written as
\[
d=r^e
\]
with exponents \(e\) lying in a union of \(M\) translates of \(\mathbb Z\),
\[
E\subseteq \bigcup_{\nu=1}^{M} (\gamma_\nu+\mathbb Z),
\qquad 0\leq \gamma_\nu<1,
\]
then from a fixed starting divisor \(d_i=r^{e_i}\),

\begin{itemize}
\item there are at most \(M-1\) later divisors whose exponent difference from \(e_i\) lies in \((0,1)\), and
\item for each integer \(k\geq 1\), there are at most \(M\) later divisors whose exponent difference lies in \([k,k+1)\).
\end{itemize}

This lets one split the tail from \(d_i\) into:

\begin{enumerate}
\item a \emph{short-range part}, controlled by the first \(M-1\) forward neighbours of \(d_i\), and
\item a \emph{long-range part}, controlled by a convergent geometric series times \(1/d_i\).
\end{enumerate}

A second new lemma then bounds the total contribution of all \(u\)-step neighbours,
\[
\sum_{i=1}^{t-u}\frac{1}{d_{i+u}-d_i},
\]
by \(\frac1u S_{\mathrm{gap}}(n)\). Summing this for \(u=1,\dots,M-1\) yields the factor \(H_{M-1}\).

This is the exact extension of the previous ``two progressions'' argument to an arbitrary fixed number \(M\) of progressions.

\section{ATTACK PLAN (ROUND-2)}

The exact gap left after Round-(round\_no-1) was this:

\begin{quote}
The previous unconditional proof handled only the case of \emph{two} exponent progressions (equivalently \(n=r^a s\)). It broke as soon as one had more than one exceptional subunit exponent difference from a given starting divisor.
\end{quote}

To overcome that, I now prove:

\begin{enumerate}
\item a general \(u\)-step gap lemma for arbitrary increasing sequences;
\item a tail lemma for exponent sets lying in a union of \(M\) translates of \(\mathbb Z\);
\item the resulting theorem for \(n=r^a m\) with \(\tau(m)=M\);
\item a sharper corollary for \(n=p^a q^b\).
\end{enumerate}

This strictly advances the previous round because it replaces the ``two progressions only'' phenomenon by a full theorem for \emph{any fixed number \(M\)} of progressions.

\section{WORK (ROUND-2)}

\subsection*{1. A general \(u\)-step gap estimate}

Let
\[
x_1<x_2<\cdots<x_t
\]
be any strictly increasing real sequence. Set
\[
g_i:=x_{i+1}-x_i \qquad (1\leq i<t).
\]

\paragraph{Lemma 1.}
For every integer \(u\) with \(1\leq u<t\),
\[
\sum_{i=1}^{t-u}\frac{1}{x_{i+u}-x_i}
\leq
\frac1u \sum_{i=1}^{t-1}\frac{1}{g_i}.
\]

\emph{Proof.}
Fix \(u\). For each \(i\in\{1,\dots,t-u\}\),
\[
x_{i+u}-x_i=g_i+g_{i+1}+\cdots+g_{i+u-1}.
\]
By Cauchy--Schwarz,
\[
\bigl(g_i+\cdots+g_{i+u-1}\bigr)\left(\frac1{g_i}+\cdots+\frac1{g_{i+u-1}}\right)\geq u^2.
\]
Hence
\[
\frac{1}{x_{i+u}-x_i}
\leq
\frac1{u^2}\sum_{k=i}^{i+u-1}\frac1{g_k}.
\]
Summing over \(i=1,\dots,t-u\) gives
\[
\sum_{i=1}^{t-u}\frac{1}{x_{i+u}-x_i}
\leq
\frac1{u^2}\sum_{i=1}^{t-u}\sum_{k=i}^{i+u-1}\frac1{g_k}.
\]
For a fixed \(k\), the term \(1/g_k\) can occur only when \(i\leq k\leq i+u-1\), that is,
\[
k-u+1\leq i\leq k.
\]
There are at most \(u\) such integers \(i\). Therefore
\[
\sum_{i=1}^{t-u}\sum_{k=i}^{i+u-1}\frac1{g_k}
\leq
u\sum_{k=1}^{t-1}\frac1{g_k}.
\]
Substituting this into the previous estimate yields
\[
\sum_{i=1}^{t-u}\frac{1}{x_{i+u}-x_i}
\leq
\frac1u\sum_{k=1}^{t-1}\frac1{g_k}.
\]
This is exactly the claim.

\subsection*{2. The \(M\)-progressions tail lemma}

Let \(r\geq 2\) be an integer. Let
\[
e_1<e_2<\cdots<e_t
\]
be real numbers, and define
\[
x_i:=r^{e_i}\qquad (1\leq i\leq t).
\]
Assume that there exist distinct real numbers
\[
0\leq \gamma_1,\dots,\gamma_M<1
\]
such that
\[
\{e_1,\dots,e_t\}\subseteq \bigcup_{\nu=1}^{M}(\gamma_\nu+\mathbb Z).
\]

Also set
\[
C_r:=\sum_{k=1}^{\infty}\frac{1}{r^k-1}.
\]

\paragraph{Lemma 2.}
For each \(i\in\{1,\dots,t-1\}\),
\[
\sum_{j=i+1}^{t}\frac{1}{x_j-x_i}
\leq
\sum_{u=1}^{\min(M-1,t-i)}\frac{1}{x_{i+u}-x_i}
+\frac{M\,C_r}{x_i}.
\]

\emph{Proof.}
Fix \(i<t\). Since \(x_j=r^{e_j}\),
\[
\sum_{j=i+1}^{t}\frac{1}{x_j-x_i}
=
\frac{1}{x_i}\sum_{j=i+1}^{t}\frac{1}{r^{e_j-e_i}-1}.
\tag{2.1}
\]

Let
\[
D_i:=\{e_j-e_i:\ j=i+1,\dots,t\}.
\]
Then every element of \(D_i\) is positive.

Write \(e_i\in \gamma_\mu+\mathbb Z\) for some \(\mu\in\{1,\dots,M\}\). If \(e_j\in \gamma_\nu+\mathbb Z\), then
\[
e_j-e_i\in (\gamma_\nu-\gamma_\mu)+\mathbb Z.
\]
Thus, for each fixed \(\nu\), the differences contributed by the \(\nu\)-th translate form a subset of a single translate of \(\mathbb Z\).

From this we obtain two counting facts.

\medskip

\noindent\textbf{Fact 1.} The interval \((0,1)\) contains at most \(M-1\) elements of \(D_i\).

Indeed, when \(\nu=\mu\), the differences are integers and hence contribute nothing to \((0,1)\). For each \(\nu\neq \mu\), the set \((\gamma_\nu-\gamma_\mu)+\mathbb Z\) contains at most one element of \((0,1)\). Summing over the \(M-1\) choices of \(\nu\neq \mu\) proves the claim.

\medskip

\noindent\textbf{Fact 2.} For every integer \(k\geq 1\), the interval \([k,k+1)\) contains at most \(M\) elements of \(D_i\).

Indeed, for each \(\nu\), the set \((\gamma_\nu-\gamma_\mu)+\mathbb Z\) contains at most one point in any interval of length \(1\). Summing over \(\nu=1,\dots,M\) gives the claim.

\medskip

Now order the elements of \(D_i\) increasingly. Since \(\Delta_u:=e_{i+u}-e_i\) are exactly the first positive differences from \(e_i\), the first
\[
\min(M-1,t-i)
\]
elements of \(D_i\) are
\[
e_{i+1}-e_i,\ e_{i+2}-e_i,\ \dots,\ e_{i+\min(M-1,t-i)}-e_i.
\]
By Fact 1, after removing these first \(\min(M-1,t-i)\) elements, every remaining element of \(D_i\) lies in \([1,\infty)\). By Fact 2, for each \(k\geq 1\), there are at most \(M\) remaining elements in \([k,k+1)\).

Since the function
\[
x\mapsto \frac{1}{r^x-1}
\]
is strictly decreasing on \((0,\infty)\), we conclude that
\[
\sum_{x\in D_i}\frac{1}{r^x-1}
\leq
\sum_{u=1}^{\min(M-1,t-i)}\frac{1}{r^{e_{i+u}-e_i}-1}
+
M\sum_{k=1}^{\infty}\frac{1}{r^k-1}.
\]
Using the definition of \(C_r\), this becomes
\[
\sum_{x\in D_i}\frac{1}{r^x-1}
\leq
\sum_{u=1}^{\min(M-1,t-i)}\frac{1}{r^{e_{i+u}-e_i}-1}
+
M C_r.
\]
Multiplying by \(1/x_i\) and recalling (2.1) yields
\[
\sum_{j=i+1}^{t}\frac{1}{x_j-x_i}
\leq
\sum_{u=1}^{\min(M-1,t-i)}\frac{1}{x_{i+u}-x_i}
+
\frac{M C_r}{x_i},
\]
which is exactly the desired bound.

\subsection*{3. A general theorem for \(n=r^a m\)}

\paragraph{Theorem 3.}
Let \(r\) be a prime, let \(a\geq 0\), and let \(m\geq 1\) satisfy \((r,m)=1\). Set
\[
M:=\tau(m).
\]
Let
\[
1=d_1<\cdots<d_t=n
\]
be the divisors of
\[
n=r^a m
\]
in increasing order. Then
\[
S_{\mathrm{all}}(n)\leq H_{M-1}S_{\mathrm{gap}}(n)+\frac{10}{3}M^2.
\]

\emph{Proof.}
List the positive divisors of \(m\) as
\[
c_1,\dots,c_M.
\]
Every divisor of \(n=r^a m\) is uniquely of the form
\[
r^u c_\nu
\qquad (0\leq u\leq a,\ 1\leq \nu\leq M).
\]
Define
\[
\gamma_\nu:=\{\log_r c_\nu\}\in [0,1).
\]
I claim that the \(\gamma_\nu\) are distinct.

Indeed, if \(\gamma_\nu=\gamma_{\nu'}\), then
\[
\log_r c_\nu-\log_r c_{\nu'}\in \mathbb Z,
\]
so
\[
\frac{c_\nu}{c_{\nu'}}=r^k
\]
for some integer \(k\). But \(c_\nu,c_{\nu'}\mid m\) and \((r,m)=1\), so the only possible power of \(r\) dividing \(m\) is \(r^0=1\). Hence \(c_\nu=c_{\nu'}\), so \(\nu=\nu'\). This proves the claim.

Thus the divisors of \(n\) can be written as
\[
r^e
\]
with exponents \(e\) lying in the union of the \(M\) translates
\[
\gamma_\nu+\mathbb Z,
\qquad 1\leq \nu\leq M.
\]
Apply Lemma 2 to the ordered divisor sequence \(d_i=r^{e_i}\), and sum over \(i=1,\dots,t-1\). This gives
\[
S_{\mathrm{all}}(n)
\leq
\sum_{u=1}^{M-1}\sum_{i=1}^{t-u}\frac{1}{d_{i+u}-d_i}
+
M C_r\sum_{i=1}^{t-1}\frac{1}{d_i}.
\tag{3.1}
\]
By Lemma 1, for each \(u\in\{1,\dots,M-1\}\),
\[
\sum_{i=1}^{t-u}\frac{1}{d_{i+u}-d_i}
\leq
\frac1u \sum_{i=1}^{t-1}\frac{1}{d_{i+1}-d_i}
=
\frac1u\,S_{\mathrm{gap}}(n).
\]
Summing over \(u\) yields
\[
\sum_{u=1}^{M-1}\sum_{i=1}^{t-u}\frac{1}{d_{i+u}-d_i}
\leq
H_{M-1} S_{\mathrm{gap}}(n).
\tag{3.2}
\]

Next,
\[
\sum_{d\mid n}\frac1d
=
\left(\sum_{u=0}^{a}\frac{1}{r^u}\right)\left(\sum_{c\mid m}\frac1c\right).
\]
Since \(r\geq 2\),
\[
\sum_{u=0}^{a}\frac1{r^u}\leq \sum_{u=0}^{\infty}\frac1{2^u}=2.
\]
Also every divisor \(c\mid m\) is at least \(1\), so
\[
\sum_{c\mid m}\frac1c \leq \tau(m)=M.
\]
Therefore
\[
\sum_{d\mid n}\frac1d \leq 2M.
\tag{3.3}
\]
In particular,
\[
\sum_{i=1}^{t-1}\frac1{d_i}\leq 2M.
\tag{3.4}
\]

Finally, because \(r\geq 2\),
\[
C_r=\sum_{k=1}^{\infty}\frac{1}{r^k-1}
\leq
\sum_{k=1}^{\infty}\frac{1}{2^k-1}.
\]
For \(k\geq 2\),
\[
2^k-1\geq 3\cdot 2^{k-2},
\]
hence
\[
\sum_{k=1}^{\infty}\frac{1}{2^k-1}
=
1+\sum_{k=2}^{\infty}\frac{1}{2^k-1}
\leq
1+\sum_{k=2}^{\infty}\frac{1}{3\cdot 2^{k-2}}
=
1+\frac23
=
\frac53.
\]
So
\[
C_r\leq \frac53.
\tag{3.5}
\]

Substituting (3.2), (3.4), and (3.5) into (3.1) gives
\[
S_{\mathrm{all}}(n)
\leq
H_{M-1}S_{\mathrm{gap}}(n)+M\cdot \frac53\cdot 2M
=
H_{M-1}S_{\mathrm{gap}}(n)+\frac{10}{3}M^2.
\]
This is exactly the required estimate.

\subsection*{4. A sharper corollary for two-prime powers}

\paragraph{Corollary 4.}
Let \(n=p^a q^b\), where \(p\neq q\) are primes and \(a,b\geq 1\). Set
\[
m:=\min(a,b).
\]
Then
\[
S_{\mathrm{all}}(n)\leq H_m\,S_{\mathrm{gap}}(n)+5(m+1).
\]
Consequently,
\[
S_{\mathrm{all}}(n)\leq \bigl(H_m+5(m+1)\bigr)\bigl(1+S_{\mathrm{gap}}(n)\bigr).
\]

\emph{Proof.}
By symmetry in \(p,q\), we may assume
\[
b\leq a,
\]
so \(m=b\). Apply Theorem 3 with
\[
r:=p,
\qquad
m_{\mathrm{cof}}:=q^b.
\]
Then
\[
\tau(m_{\mathrm{cof}})=b+1=m+1.
\]
Thus Lemma 2 and the proof of Theorem 3 give
\[
S_{\mathrm{all}}(n)
\leq
H_m\,S_{\mathrm{gap}}(n)+(m+1)C_p\sum_{d\mid n}\frac1d.
\tag{4.1}
\]

It remains to bound \(\sum_{d\mid n}1/d\). Since the divisors are \(p^u q^v\) with \(0\leq u\leq a\), \(0\leq v\leq b\),
\[
\sum_{d\mid n}\frac1d
=
\left(\sum_{u=0}^{a}\frac1{p^u}\right)\left(\sum_{v=0}^{b}\frac1{q^v}\right)
\leq
\frac{1}{1-\frac1p}\cdot \frac{1}{1-\frac1q}.
\]
Because \(p,q\) are distinct primes, the largest possible value of the right-hand side occurs for \(\{p,q\}=\{2,3\}\), giving
\[
\sum_{d\mid n}\frac1d \leq 2\cdot \frac32 = 3.
\tag{4.2}
\]
Also \(C_p\leq \frac53\) by the proof of Theorem 3. Substituting this and (4.2) into (4.1) yields
\[
S_{\mathrm{all}}(n)
\leq
H_m\,S_{\mathrm{gap}}(n)+(m+1)\cdot \frac53\cdot 3
=
H_m\,S_{\mathrm{gap}}(n)+5(m+1).
\]
This proves the first claim. The second follows immediately from
\[
H_m\,S_{\mathrm{gap}}(n)+5(m+1)\leq \bigl(H_m+5(m+1)\bigr)\bigl(1+S_{\mathrm{gap}}(n)\bigr).
\]
This proves Corollary 4.

\subsection*{5. Consequences}

For every fixed integer \(B\geq 1\), if
\[
n=p^a q^b
\qquad\text{with}\qquad
\min(a,b)\leq B,
\]
then Corollary 4 gives
\[
S_{\mathrm{all}}(n)\leq C_B\bigl(1+S_{\mathrm{gap}}(n)\bigr),
\qquad
C_B:=H_B+5(B+1).
\]

Thus the conjectured type of estimate is now proved unconditionally for all two-prime numbers with one exponent bounded, and more generally for all numbers \(n=r^a m\) with \(\tau(m)\) bounded.

\section{ADVERSARIAL VERIFICATION}

I now check the new work against the main failure modes.

\paragraph{1. Distinct fractional parts in Theorem 3.}
The proof used that if
\[
\{\log_r c_\nu\}=\{\log_r c_{\nu'}\},
\]
then \(c_\nu/c_{\nu'}=r^k\) for some integer \(k\). Because both \(c_\nu,c_{\nu'}\mid m\) and \((r,m)=1\), the quotient cannot be a nontrivial power of \(r\). Hence \(k=0\) and \(c_\nu=c_{\nu'}\). This step is valid.

\paragraph{2. Counting of short differences in Lemma 2.}
For a fixed translate \((\gamma_\nu-\gamma_\mu)+\mathbb Z\), an interval of length \(1\) contains at most one point. Therefore \((0,1)\) contains at most one point from each translate, and \([k,k+1)\) contains at most one point from each translate. Summing over \(M\) translates gives the stated bounds. No hidden multiplicity was missed.

\paragraph{3. End effects near \(i=t\).}
Lemma 2 uses
\[
\min(M-1,t-i)
\]
rather than \(M-1\). Thus if there are fewer than \(M-1\) later divisors, the statement still makes sense and is correct.

\paragraph{4. The \(u\)-step gap lemma.}
In Lemma 1, for fixed \(k\), the term \(1/g_k\) appears for at most \(u\) values of \(i\). This is the only combinatorial input in that proof, and it is correct.

\paragraph{5. Conversion to a \(C_M(1+S_{\mathrm{gap}})\) bound.}
From Theorem 3 we have
\[
S_{\mathrm{all}}(n)\leq H_{M-1}S_{\mathrm{gap}}(n)+\frac{10}{3}M^2.
\]
Since
\[
H_{M-1}S_{\mathrm{gap}}(n)\leq C_M(1+S_{\mathrm{gap}}(n))
\quad\text{and}\quad
\frac{10}{3}M^2\leq C_M(1+S_{\mathrm{gap}}(n))
\]
for
\[
C_M:=H_{M-1}+\frac{10}{3}M^2,
\]
the corollary is valid.

\paragraph{6. Interaction with the previous round.}
There is no contradiction with the earlier conditional disproof. The new theorem applies only when the divisor exponents lie in a \emph{bounded} number \(M\) of translates of \(\mathbb Z\). Tao's conditional counterexamples come from products of many nearly consecutive primes, where the number of such translates grows with the construction parameter.

\paragraph{7. Sharp remaining obstruction.}
The present method makes the dependence on \(M\) explicit. The first unresolved barrier is therefore now sharply isolated:

\begin{quote}
Can one control the contribution of the first \(M-1\) exceptional short exponent differences \emph{uniformly in \(M\)}?
\end{quote}

Without a new idea, Lemma 2 necessarily produces an \(M\)-dependent term, and this is exactly where the proof stops short of the full conjecture.

\section{FINAL}

\textbf{UNRESOLVED (BUT STRICTLY ADVANCED).}

This round proves a theorem strictly stronger than the previous round:

\begin{quote}
If \(n=r^a m\) with \(r\) prime, \((r,m)=1\), and \(\tau(m)\leq M\), then
\[
S_{\mathrm{all}}(n)\leq H_{M-1}S_{\mathrm{gap}}(n)+\frac{10}{3}M^2,
\]
hence
\[
S_{\mathrm{all}}(n)\leq C_M\bigl(1+S_{\mathrm{gap}}(n)\bigr)
\]
for an explicit constant \(C_M\) depending only on \(M\).
\end{quote}

In particular, for \(n=p^a q^b\) with
\[
m:=\min(a,b),
\]
one has the sharper bound
\[
S_{\mathrm{all}}(n)\leq H_m\,S_{\mathrm{gap}}(n)+5(m+1).
\]

So the original Erd\H{o}s problem remains open unconditionally, but the unconditional positive family has been enlarged from the previous two-progression case to all fixed-\(M\) progression families of the form \(r^a m\) with \(\tau(m)\) bounded.

\section{COMPLETION ESTIMATE}

\noindent \textbf{COMPLETION: 88\%.}

\section{REFERENCES}

\begin{enumerate}
\item T. Tao, \emph{On the sum of reciprocals of gaps between divisors}, September 9, 2025.
\end{enumerate}
